% Beispiel für das Corporate Design der TU Chemnitz (Stand 2019)
\DocumentMetadata{pdfversion=1.7}
\documentclass{ltx-talk}

\usepackage[fakcolor=if]{ltxtalk-theme-tuc-2019}
\settucthemelogo{../../images/logos/tuc_white.pdf}
%\settucthemetitlelogo{../../images/abstrakt.jpg}
\settucthemeurl{https://www.tu-chemnitz.de/}
\useltxtalktheme{tuc-2019}

\title{Echtzeitsysteme}
\subtitle{0. Kapitel\\Organisatorisches}
\author{Prof. Matthias Werner}
\institute{Professur Betriebssysteme}
\date{Sommersemester 2026}

\begin{document}
\settucthemeheader{teaching}

\maketuctitle

\section{Organisatorisches}

\subsection{Willkommen}


\begin{frame}
  \frametitle{Willkommen zur Lehrveranstaltung}
  \framesubtitle{Echtzeitsysteme}

  Das Theme bildet die wesentlichen Elemente des TU-Chemnitz-Designs ab:
  \begin{itemize}
    \item grünes Kopfband mit Universitätslogo und Veranstaltungskontext,
    \item grüne Folientitel und Aufzählungszeichen,
    \item dreigeteilte Fußzeile mit Autor, Foliennummer und Webadresse.
  \end{itemize}
\end{frame}

\settucthemeheader{section}
\tucthreeheader
\begin{frame}
  \frametitle{Dreizeilige Kopfzeile}
  \begin{itemize}
    \item Präsentationstitel
    \item aktuelle Section
    \item aktuelle Subsection
  \end{itemize}
\end{frame}

\settucthemeleftlogos{../../images/logos/osg.png,../../images/logos/tuckhs_color.pdf}
\tucnarrowframe
\begin{frame}
  \frametitle{Schmaler Inhaltsbereich}
  Mehrere Logos werden untereinander auf die Breite der linken Fläche
  skaliert. Der Inhalt beginnt auf der Fluchtlinie des Titelblocks.
\end{frame}

\tucwideframe
\tuctwoheader
\settucthemeheader{title}
\maketitle[talk]

\settucthemetitleimage[center]{../../images/web.png}
\maketitle[image]

\end{document}
